% ======================= CHAPTER VI =======================
\chapter{Conclusions and Future Work}\label{ch:discussion-and-conclusions}

\section{Future Work}

There are several areas for future work to enhance the performance and reliability of the MAGIC system. These areas includes:

\begin{itemize}
    \item \textbf{Adaptive Thresholding and Detection}
    \item \textbf{Channel Estimation / Tracking Approaches}
    \item \textbf{Time/Frequency Synchronization Enhancements}
    \item \textbf{Intelligent Tunning}
    \item \textbf{Optimize Coupling Variability and Misalignment}
\end{itemize}
%\subsection{Filtering and Noise Reduction}
% TODO: Narrowband filters, adaptive filtering.

%\subsection{Adaptive Thresholding and Detection}\label{subsec:adaptive-detection}
% TODO: Non-coherent detection improvements.

%\subsection{Channel Estimation / Tracking Approaches}\label{subsec:channel-estimation}
% TODO: Estimating effective gain via pilots.

%\subsection{Time/Frequency Synchronization Enhancements}\label{subsec:dsp-synchronization}
% TODO: Sliding correlation.

%\subsection{Frequency Drift and Detuning}\label{subsec:frequency-drift}
%\subsection{Coupling Variability and Misalignment}\label{subsec:coupling-variability}
%\subsection{Interference and Regulatory Considerations}\label{subsec:interference}

%\subsection{Bandwidth and Attenuation}\label{subsec:attenuation-bandwidth}
%\subsection{Coil Size vs Portability Constraints}\label{subsec:coil-size}

%\subsection{Signal Processing (DSP) Techniques in TTE MI Systems}\label{sec:dsp-tte}
% (From: DSP in TTE)

%\subsection{Robustness to Misalignment and Orientation}\label{subsec:robustness-misalignment}
% TODO: Orientation statistics.

\section{Conclusions}
In conclusion the developement of MAGIC has demonstrated the feasibility of a complete Through the Earth (TTE) communication system for underground environments using relatively simple hardware and software components. The system has shown promising results in achieving text messaging through at least 200 meters rock and solid material. MAGIC is a valuable tool for mining and other industries that need a reliable communication solution in challenging underground environments, considering relatively low cost.
\newp
The iterative design and testing process allowed the optimization of the system's performance, finally reaching a 100\% success rate in the communication tests considering detection and SER (Symbol Error Rate) lower tan 5\%. The use of Reed Solomon error correction with 20 parity symbols permitted to reduce the SER from 5.57\% to 4.48\%, so consists not a great improvement, but it could be useful if the encoding model change to a symbol per character as it was explained in section \ref{sec: Performance evaluation at Cerro Calan}.
\newp
The model of preamble with 24 symbols demonstrated lower variability in the correlation peak compared to the model with 12 symbols, enhancing the reliability of preamble detection and leading to more consistent performance in message reception. But the cost of processing the correlation by parts algorithm is higher, that's why the model with 12 symbols was used in the final implementation of MAGIC, and is the current model implemented in Raspberry Pi. However, the model with 24 symbols could be implemented in future works if the processing power of the hardware is improved.
\newp
In relation of the performance of the system in operational environments it's possible to conclude that the system works properly at Benjamin's mine, specially considering detection. However it's clear that the demodulation performance is strongly affected by the noise spectral distribution in the mine, which is more impulsive and with higher amplitude compared to the noise in the astronomical observatory. The type of interference that's is sudden and short in time particularly affects the performance of the demodulator due to max peak detection. The next step is know better the characteristics of the noise in the mine and other environments to adapt the system and subtract the mean noise spectral distribution to improve SER.
